\newpage
\section*{АНОТАЦІЯ}
\addcontentsline{toc}{section}{АНОТАЦІЯ}

У курсовій роботі розроблено програмну систему діалогової взаємодії з великими мовними моделями. Система призначена для створення та ведення чатів з локальною LLM, збереження історії повідомлень, пошуку по діалогах, експорту чату у Markdown, перемикання локальних моделей і використання LoRA/QLoRA-адаптерів.

Програмний продукт реалізовано мовою Python. Для web-інтерфейсу використано Flask, для збереження даних --- PostgreSQL, для керування міграціями --- Alembic, для запуску локальної мовної моделі --- Hugging Face Transformers. Додатково реалізовано CLI-режим, потокову генерацію відповідей, кнопку зупинки генерації, generation presets, локальний model registry, підтримку математичних формул через KaTeX і підсвічування блоків коду.

У роботі сформовано функціональні та нефункціональні вимоги до системи, виконано аналіз предметної області, спроєктовано багаторівневу архітектуру, ER-діаграму бази даних, UML-діаграми прецедентів, класів і послідовності. Під час реалізації використано принципи ООП, SOLID, розділення відповідальності, Repository/Storage Adapter, Service Layer та Adapter.

Працездатність системи перевірено автоматизованими тестами. Повний набір містить 105 успішних тестів, а покриття production-коду, отримане за допомогою \texttt{coverage.py}, становить 77\%.

\textbf{Ключові слова:} велика мовна модель, LLM, Flask, PostgreSQL, Transformers, LoRA, QLoRA, web-застосунок, CLI, streaming generation, тестування.

